\chapter{Introduzione del modello}

Introduciamo un modello matematico a tempo discreto con l'obiettivo di rappresentare in modo sintetico come la composizione del gruppo e lo stato di affaticamento medio varino nel tempo in risposta al carico di allenamento, agli infortuni e ai flussi di ingresso e uscita dalla popolazione. Il modello non ha la pretesa di descrivere il comportamento individuale di ogni atleta, ma di catturare le tendenze aggregate di un gruppo sufficientemente numeroso, per il quale le dinamiche collettive sono più rilevanti delle fluttuazioni individuali. In altre parole, il modello si colloca a un livello di astrazione intermedio, dato che è sufficientemente dettagliato da distinguere stati diversi degli atleti e da descrivere l'evoluzione della fatica, ma sufficientemente aggregato da rimanere trattabile analiticamente e da non richiedere la simulazione di ogni singolo individuo. \par

Il modello è formulato a tempo discreto, con $t = 0, 1, \dots, T$, dove $T$ è l'orizzonte temporale considerato. Si assume che le sedute di allenamento siano distribuite circa omogeneamente nel tempo, così che ogni passo $t \to t+1$ corrisponda a una seduta e il tempo che intercorre tra due sedute sia approssimativamente lo stesso. \\
Se consideriamo ad esempio un programma di allenamento basato su una seduta al giorno, possiamo dire che ogni passo $t$ corrisponda a un giorno. In questo modo è possibile confrontare anche il recupero tra due sedute consecutive, perché l'intervallo di tempo tra un passo e il successivo è pressoché costante per ipotesi. \par

Questa scelta è motivata dal fatto che gli allenamenti sono eventi discreti, separati da periodi di riposo, e che il modello intende descrivere l'evoluzione del gruppo seduta dopo seduta. L'ipotesi di distribuzione omogenea permette inoltre di trattare il tempo come una variabile regolare e di evitare la complessità aggiuntiva legata a pianificazioni irregolari, senza perdere la capacità di descrivere i fenomeni principali. L'ipotesi di equidistanza non è restrittiva nella pratica, dato che un eventuale giorno di riposo può essere trattato come una seduta a carico nullo, senza modificare la struttura del modello.

In questa prima formulazione, il modello non include un parametro di controllo. Le funzioni che descrivono l'efficacia della prevenzione e i carichi di allenamento sono trattate come parametri esogeni, assegnati dall'esterno. Questa scelta permette di studiare il comportamento del sistema in risposta a diversi scenari di allenamento e prevenzione, senza aggiungere la complessità dell'ottimizzazione del parametro. Pertanto l'aggiunta del controllo è rinviata a una fase successiva. In questo modo è possibile analizzare separatamente due aspetti distinti, cioè da un lato la dinamica del sistema in condizioni assegnate, dall'altro la ricerca delle condizioni ottimali. Questa separazione è comune in modellistica e permette di costruire una comprensione solida del sistema prima di affrontare il problema più complesso dell'ottimizzazione.

\section{Variabili}

Gli atleti sono suddivisi in quattro popolazioni, che rappresentano stati diversi rispetto alla condizione fisica e alla capacità di allenarsi. La suddivisione riflette la necessità di distinguere non solo tra atleti sani e infortunati, ma anche tra diversi livelli di rischio e capacità di allenarsi. Questa distinzione è importante perché il rischio di infortunio, la capacità di partecipare all'allenamento e la dinamica di recupero non sono uniformi tra tutti gli atleti, e trattarli come un unico gruppo omogeneo comporterebbe una perdita significativa di informazione.

Le popolazioni sono:
\begin{itemize}
    \item Atleti sani $S$: atleti non infortunati e senza infortuni pregressi. Si possono allenare senza limitazioni e sono considerati il gruppo di riferimento per il rischio di infortunio. La loro condizione fisica è ottimale e non presentano vulnerabilità residue legate a eventi precedenti.
    \item Atleti con pregressi $P$: atleti che hanno avuto infortuni prima dell'inizio del periodo preso in considerazione. Pur essendo pienamente allenabili, sono considerati più a rischio rispetto ai sani, in quanto un infortunio pregresso può lasciare una vulnerabilità residua, una ridotta stabilità articolare o una minore tolleranza al carico. Questi atleti partecipano all'allenamento senza restrizioni formali, ma la loro storia clinica giustifica un trattamento differenziato nel modello del rischio.
    \item Atleti infortunati $I$: atleti che si sono infortunati durante il periodo preso in considerazione e non si stanno allenando. Sono temporaneamente esclusi dall'attività e non contribuiscono alla densità sul tatami. La loro presenza è rilevante perché rappresenta l'esito più grave del rischio di infortunio e perché il loro recupero parziale alimenta la popolazione $M$.
    \item Atleti ``malconci'' $M$: atleti con infortuni lievi o che sono parzialmente guariti dall'infortunio e possono continuare ad allenarsi con qualche restrizione. Partecipano all'allenamento, ma la loro condizione fisica non è ottimale e il loro rischio di incorrere in un nuovo infortunio è considerato maggiore rispetto a quello degli atleti sani. Questi atleti rappresentano una categoria intermedia: non sono abbastanza compromessi da dover sospendere l'attività, ma non sono nemmeno in condizioni ottimali.
\end{itemize}

La scelta di distinguere quattro popolazioni, invece di due o tre, è motivata dalla necessità di rappresentare separatamente tre aspetti distinti, cioè il rischio di infortunio, che è diverso per sani, pregressi e malconci, la capacità di allenarsi, che distingue gli infortunati dagli altri, e il recupero parziale, che non riporta l'atleta allo stato di sano ma lo rende nuovamente allenabile. Se avessimo distinto solo sani e infortunati, non avremmo potuto differenziare il rischio tra sani e pregressi, né modellare il recupero parziale. Se avessimo distinto solo tre popolazioni, avremmo dovuto rinunciare a una di queste distinzioni, con una perdita di realismo. \par

Le popolazioni $S$, $P$, $I$, $M$ sono rappresentate come numeri reali non negativi, non come conteggi interi. Questa scelta è giustificata dal fatto che il modello è pensato per gruppi sufficientemente numerosi, per i quali le fluttuazioni sono trascurabili e l'approssimazione continua è adeguata. Inoltre, la rappresentazione continua permette di utilizzare strumenti analitici che sarebbero più complessi in un modello a conteggi interi, come l'analisi di stabilità e lo studio dei punti di equilibrio.

\subsection{Accumulo di fatica}

A queste popolazioni si aggiungono due variabili di stato che descrivono l'accumulo di fatica negli atleti:
\begin{itemize}
    \item Fatica fisica del gruppo $\bar{F}_f$: rappresenta il livello medio di affaticamento muscolare e articolare del gruppo che si allena. Aumenta con il carico di allenamento e diminuisce con il recupero.
    \item Fatica psicologica del gruppo $\bar{F}_p$: rappresenta il livello medio di affaticamento mentale del gruppo che si allena. Aumenta con il carico di allenamento e con lo stress legato alla competizione, e diminuisce con il recupero.
\end{itemize}

Le due componenti della fatica sono trattate separatamente perché evolvono con dinamiche proprie e contribuiscono in modo distinto al rischio di infortunio. Un atleta può essere fisicamente fresco ma mentalmente affaticato, o viceversa. La separazione permette di modellare scenari in cui il carico fisico e quello psicologico non sono proporzionali, e di studiare l'effetto di ciascuna componente sul rischio complessivo. Inoltre, le due componenti possono avere tassi di recupero diversi, e la loro separazione permette di catturare questa differenza senza forzare una relazione di proporzionalità tra fatica fisica e psicologica. \par

La fatica è calcolata come media sul gruppo che si allena, cioè sugli atleti in $S$, $P$ e $M$. Gli atleti in $I$ non contribuiscono al calcolo, perché non si allenano e non sono esposti al carico. Questa scelta riflette il fatto che la fatica rilevante per il rischio di infortunio è quella degli atleti che effettivamente partecipano all'allenamento. Se gli atleti infortunati contribuissero alla media, la fatica media risulterebbe artificialmente diluita, perché gli infortunati sono per definizione a riposo. \par

Le variabili $\bar{F}_f$ e $\bar{F}_p$ sono rappresentate come numeri reali non negativi. La loro evoluzione è descritta da equazioni che bilanciano il carico aggiunto in ogni seduta e il recupero nel tempo. In ogni passo $t \to t+1$, il gruppo accumula una quantità di fatica proporzionale al carico della seduta, e ne smaltisce una parte proporzionale al livello di fatica corrente.

\section{Ipotesi iniziali}

Per garantire la coerenza e la non distorsione del modello, richiediamo una serie di ipotesi iniziali che non dipendono dalla struttura delle equazioni. Queste ipotesi definiscono il contesto in cui il modello è valido e delimitano le situazioni in cui le sue conclusioni possono essere applicate. \par

In primis, essendo il tempo calcolato in modo discreto $t = 0, 1, \dots, T$, si suppone che le sedute di allenamento siano distribuite equidistanti nel tempo, così che ogni passo del modello $t \to t+1$ abbia approssimativamente la stessa lunghezza. Questo viene richiesto per poter calcolare correttamente la decadenza della fatica accumulata, dato che se gli intervalli tra le sedute fossero di lunghezza diversa, il recupero tra due passi consecutivi non sarebbe confrontabile, e la dinamica della fatica risulterebbe distorta. L'ipotesi di equidistanza permette di trattare il tempo come una variabile regolare e di utilizzare le stesse funzioni di recupero per ogni passo. Un eventuale giorno di riposo, cioè una seduta di allenamento saltata, può essere considerato come una seduta a carico di fatica nullo e trattato all'interno delle variabili di accumulo della fatica, senza bisogno di modificare la struttura del modello. \par

Per poter considerare gli atleti come gruppo e non come singoli individui, sia in termini di accumulo di fatica sia di probabilità di infortunio, supponiamo che abbiano una preparazione fisica e tecnica omogenea, quindi che non ci siano differenze marcate tra di loro. Questa ipotesi è necessaria per giustificare l'uso di variabili medie, cioè $\bar{F}_f$ e $\bar{F}_p$, al posto delle fatiche individuali e per poter applicare lo stesso tasso di infortunio a tutti gli atleti di una data popolazione. Se il gruppo fosse eterogeneo, la media non sarebbe rappresentativa e il modello perderebbe di significato. L'omogeneità non richiede che gli atleti siano identici, ma che le differenze siano trascurabili rispetto alla scala dei fenomeni modellati. \par

L'orizzonte temporale è breve, ad esempio una settimana, un mese o una stagione sportiva. Questo fa sì che gli atleti infortunati $I$ o malconci $M$ non abbiano sufficiente tempo di recuperare completamente e tornare nella popolazione $P$ o $S$. Di conseguenza, nel modello è previsto un solo percorso di recupero, quello da infortunati $I$ a malconci $M$, che rappresenta una guarigione parziale. Il ritorno completo allo stato di sano o di pregresso è escluso perché richiederebbe tempi più lunghi di quelli considerati. Questa ipotesi semplifica la struttura del modello e riflette la realtà di un periodo di allenamento, in cui gli infortuni gravi raramente si risolvono completamente prima della fine del periodo. \par

Le popolazioni sono numeri reali, non conteggi interi. Il modello quindi è significativo e non distorto se il numero di atleti è sufficientemente numeroso. Questa ipotesi permette di utilizzare strumenti analitici continui e di evitare le complicazioni legate alla discretizzazione dei conteggi. Per gruppi piccoli, le fluttuazioni potrebbero diventare rilevanti e l'approssimazione continua potrebbe perdere accuratezza, ma per gruppi sufficientemente numerosi le dinamiche aggregate sono ben descritte dal modello continuo. \par

A queste ipotesi si aggiungono quelle discusse nelle sezioni successive, riguardanti le equazioni del modello e le funzioni usate. In particolare, le ipotesi sul mescolamento degli accoppiamenti, sulla simmetria del rischio di infortunio e sulla numerosità del gruppo sono necessarie per giustificare l'uso delle fatiche medie al posto delle fatiche individuali e per garantire che il modello catturi in modo adeguato le dinamiche collettive.

\section{Condizioni iniziali}

Le condizioni iniziali del modello sono:
\begin{itemize}
    \item $S(0) \geq 0$: il gruppo può iniziare con un numero arbitrario di atleti sani.
    \item $P(0) \geq 0$: il gruppo può iniziare con un numero arbitrario di atleti con infortuni pregressi.
    \item $I(0) = 0$: si assume che nessun atleta sia infortunato all'inizio del periodo considerato.
    \item $M(0) = 0$: si assume che nessun atleta sia malconcio all'inizio del periodo considerato.
    \item $\bar{F}_f(0) \geq 0$: la fatica fisica può essere presente all'inizio del periodo considerato, ma si preferisce considerarla nulla.
    \item $\bar{F}_p(0) \geq 0$: la fatica psicologica può essere presente all'inizio del periodo considerato, ma si preferisce considerarla nulla.
\end{itemize}
Si deve avere inoltre che $S(0) + P(0) > 0$, cioè che la popolazione iniziale non sia vuota. Senza atleti sani o con infortuni pregressi, il modello non avrebbe una popolazione di partenza su cui evolvere, e le equazioni perderebbero di significato. La condizione $I(0) = M(0) = 0$ è invece una scelta semplificativa, non un requisito di coerenza. \par

La scelta $I(0) = M(0) = 0$ è coerente con l'idea che il modello descriva un periodo di allenamento che inizia in condizioni di relativa normalità, senza infortuni in corso. Se il periodo considerato iniziasse immediatamente dopo un infortunio, o in una fase di recupero, sarebbe più appropriato porre $I(0) > 0$ e $M(0) > 0$. In questa prima formulazione si assume che il gruppo parta da una condizione di salute, e che eventuali infortuni siano il risultato dell'evoluzione del modello, non una condizione iniziale.

Le condizioni iniziali su $\bar{F}_f(0)$ e $\bar{F}_p(0)$ ammettono un livello di fatica non nullo all'inizio, ma si preferisce considerarlo nullo. In questo modo si suppone che il gruppo parta da uno stato di freschezza, eventualmente dopo un periodo di riposo o all'inizio della stagione. Questa scelta semplifica il problema di dover calcolare il livello di fatica medio prima dell'arrivo, e riflette l'ipotesi che il periodo considerato inizi dopo una fase di recupero sufficiente a smaltire la fatica accumulata in precedenza. Condizioni iniziali diverse possono essere considerate per modellare scenari in cui il gruppo entra nel periodo di allenamento con fatica residua, ma in questa prima formulazione si assume che la fatica iniziale sia nulla.

La scelta di porre $\bar{F}_f(0) = \bar{F}_p(0) = 0$ non è obbligatoria, ma è conveniente per due motivi. In primo luogo, semplifica l'analisi del modello, eliminando la necessità di stimare la fatica iniziale del gruppo. In secondo luogo, rende più chiari i confronti tra scenari diversi, perché tutti partono dallo stesso stato di freschezza. Se si volesse modellare un gruppo che entra nel periodo con fatica residua, basterebbe porre $\bar{F}_f(0) > 0$ e/o $\bar{F}_p(0) > 0$, senza modificare la struttura delle equazioni.

La scelta di porre $S(0) \geq 0$ e $P(0) \geq 0$, con $S(0) + P(0) > 0$, riflette il fatto che il gruppo può iniziare con una composizione arbitraria di atleti sani e con infortuni pregressi. Il modello non richiede che tutti gli atleti siano sani, né che tutti abbiano infortuni pregressi. La composizione iniziale è un parametro del modello, che può essere variato per studiare scenari diversi.

\section{Funzioni}

Nelle equazioni del modello sono presenti cinque funzioni, che regolano i processi di ingresso, uscita, infortunio, carico e recupero. Esse costituiscono il legame tra i parametri e la dinamica delle variabili di stato, e la loro specificazione è necessaria per poter simulare e analizzare il modello.

\begin{itemize}
    \item Funzione di crescita $\mu$: descrive il numero di nuovi atleti che entrano nel gruppo complessivo in ogni seduta. Riflette le dinamiche di iscrizione a una stagione sportiva, con ingressi tipicamente concentrati all'inizio del periodo e progressivamente decrescenti. La funzione dipende dalla frazione di tempo trascorsa e da parametri che ne controllano l'ampiezza e la forma.
    \item Funzione di abbandono $\nu$: descrive la frazione di atleti che abbandona il gruppo in ogni seduta, per motivi diversi. Il suo profilo temporale può essere opposto a quello della funzione di crescita, dato che gli abbandoni tendono a concentrarsi alla fine del periodo, quando l'entusiasmo iniziale si riduce o sopraggiungono altri impegni. La funzione è modulata da un moltiplicatore specifico per la popolazione a cui si applica.
    \item Funzione di infortunio $\lambda$: descrive la probabilità che un atleta si infortuni in una seduta. Dipende dalla fatica fisica e psicologica media del gruppo, dalla densità di atleti sul tatami e da un moltiplicatore di rischio specifico per la popolazione a cui è applicata. La funzione è crescente rispetto a fatica e densità e si satura a un valore massimo.
    \item Funzione di carico $c$: fornisce il tasso di carico di ogni seduta di allenamento per le due componenti della fatica, fisica e psicologica. Può essere costante, periodica o variabile nel tempo, e rappresenta l'entrata principale che determina l'accumulo di fatica.
    \item Funzione di recupero $r$: descrive il tasso di smaltimento della fatica in ogni seduta. Il recupero può essere massimo all'inizio del periodo, quando gli atleti sono freschi, e decrescere nel tempo, man mano che la fatica si accumula.
\end{itemize}

Queste funzioni verranno specificate in dettaglio più avanti, dove verranno definite le loro forme analitiche, i parametri, i domini e le proprietà. In questa sede è sufficiente averne un'idea qualitativa, dato che $\mu$ e $\nu$ regolano i flussi di popolazione, $\lambda$ regola il rischio di infortunio, $c$ fornisce il carico di allenamento e $r$ regola la dinamica della fatica. Le loro interazioni determinano l'evoluzione del sistema e la composizione del gruppo nel tempo.